\documentclass[11pt]{article}

\usepackage[utf8]{inputenc}
\usepackage[T1]{fontenc}
\usepackage{lmodern}
\usepackage{microtype}
\usepackage{geometry}
\usepackage{hyperref}
\usepackage{booktabs}
\usepackage{listings}
\usepackage{xcolor}
\usepackage{enumitem}

\geometry{a4paper, margin=1in}
\hypersetup{
    colorlinks=true,
    linkcolor=black,
    urlcolor=blue!60!black,
    pdftitle={Social Media Ingestion for the Global Flood Event Data Pipeline}
}

\definecolor{codebg}{rgb}{0.97,0.97,0.97}
\definecolor{codecomment}{rgb}{0.35,0.35,0.35}
\definecolor{codekw}{rgb}{0.10,0.30,0.60}

\lstset{
    backgroundcolor=\color{codebg},
    basicstyle=\ttfamily\footnotesize,
    breaklines=true,
    columns=fullflexible,
    frame=single,
    framerule=0pt,
    keywordstyle=\color{codekw},
    commentstyle=\color{codecomment},
    showstringspaces=false,
    xleftmargin=8pt,
    xrightmargin=8pt,
    aboveskip=8pt,
    belowskip=8pt
}

\title{Social Media Ingestion for the Global Flood Event Data Pipeline\\
       \large A technical report on design, integration, and production cutover}
\author{Mohammed Chadli}
\date{May 2026}

\begin{document}

\maketitle

\begin{abstract}
This report documents how I added a complete social media evidence layer to
my Global Flood Event Data Pipeline. The goal was not to invent new flood
events from public posts, but to add a corroborating signal on top of the
existing authoritative sources (Dartmouth FO, GloFAS, Copernicus EMS,
EM-DAT, ReliefWeb). I built the pipeline locally over twelve incremental
steps, hit a 401/403 wall during the production cutover, debugged it down
to the wrong Bluesky host being used for authenticated reads, fixed it,
then spent two further iterations tightening the filtering and adding
text-based country extraction. The system is now live against Supabase:
50~real posts ingested through the authenticated Bluesky PDS,
46~normalised staging signals after the self-healing transformer dropped
4~political metaphors, 61\% of those signals carry an inferred country,
average signal confidence 0.77, and all eight social data-quality checks
pass with zero failing rows. The full implementation is covered by
138~automated tests.
\end{abstract}

\tableofcontents

\section{Motivation and Design Principles}

The pipeline already ingests authoritative flood records from five
official sources. Public social media data is qualitatively different: it
is timely, geographically broad, and very noisy. I wanted it in the
pipeline as an early-warning evidence layer, but never as truth on its
own.

Three principles shaped every decision:

\begin{enumerate}[leftmargin=*]
  \item \textbf{Social posts are evidence, not events.} A single post can
        never enter \texttt{staging.flood\_events}. Social signals live in
        their own staging table and are joined to events only inside a
        dedicated mart that downstream consumers can opt into.
  \item \textbf{Raw payloads are preserved verbatim.} Whatever the API
        returns is stored as JSONB, with the same audit columns
        (\texttt{source}, \texttt{batch\_id}, \texttt{ingested\_at},
        \texttt{file\_path}) used by every other ingestion module. The
        whole pipeline is therefore re-runnable from raw without going
        back to Bluesky.
  \item \textbf{Filtering must stay explainable.} Every rejection or
        admission decision is recorded inline on the row so an analyst can
        ask ``why was this post kept?'' or ``why was this one dropped?''
        without having to re-run the classifier.
\end{enumerate}

The end-to-end shape is:

\begin{lstlisting}
Bluesky search (authenticated PDS, public AppView fallback)
   |
   v
raw.social_media_posts
   |
   v
staging.social_flood_signals  (normalized + scored)
   |
   v
marts.social_flood_signals
marts.flood_events_with_social_signals
marts.social_signals_by_country_day
   |
   v
FastAPI endpoints
   |
   v
Airflow orchestration + data-quality report
\end{lstlisting}

\section{Choice of Platform}

I evaluated four candidates (Twitter/X, Reddit, Mastodon, Bluesky) and
picked Bluesky as the first integration. The reasoning:

\begin{itemize}[leftmargin=*]
  \item \textbf{Open AT Protocol.} The
        \texttt{app.bsky.feed.searchPosts} lexicon is documented, stable,
        and reachable from a plain HTTPS GET. No proprietary SDK needed.
  \item \textbf{Two access tiers.} An unauthenticated AppView at
        \texttt{public.api.bsky.app} for development and a fully
        authenticated PDS at \texttt{bsky.social}. The two-tier design
        meant I could prototype without credentials and switch to authed
        access in production with the same code path.
  \item \textbf{App passwords.} A separate credential type (App Password)
        designed for third-party clients, revocable independently of the
        account password. The pipeline never needs my real password.
  \item \textbf{Bounded search.} A simple keyword search endpoint with a
        \texttt{limit} parameter (capped at~100) and ISO~timestamp
        \texttt{since}/\texttt{until} filters. That maps cleanly onto a
        bounded Airflow task, no firehose state machine needed.
  \item \textbf{Topical fit.} Real weather services (NWS regional
        accounts, civil-protection accounts in Europe) already post
        flood advisories there. The first 50 real posts I ingested
        contained NWS Anchorage flood advisories and Spanish
        \emph{Aemet} warnings.
\end{itemize}

Mastodon, Reddit, and X/Twitter are kept on the roadmap behind a single
platform-agnostic \texttt{raw.social\_media\_posts} table so adding any
of them is a new ingester plus a normalizer, not a schema change.

\section{Data Contract and Schema}

The schema lives in \texttt{db/schema.sql}. Two new tables were added.

\subsection{\texttt{raw.social\_media\_posts}}

I deliberately picked one unified raw table over per-platform tables. A
\texttt{(platform, post\_id)} composite unique constraint guarantees
idempotency across reruns.

\begin{lstlisting}[language=SQL]
CREATE TABLE IF NOT EXISTS raw.social_media_posts (
    id              BIGSERIAL PRIMARY KEY,
    platform        TEXT NOT NULL,
    source          TEXT NOT NULL,
    source_url      TEXT,
    post_id         TEXT NOT NULL,
    file_path       TEXT,
    batch_id        TEXT NOT NULL,
    ingested_at     TIMESTAMPTZ NOT NULL DEFAULT NOW(),
    payload         JSONB NOT NULL,
    CONSTRAINT uq_social_media_posts_platform_post
        UNIQUE (platform, post_id)
);
\end{lstlisting}

\subsection{\texttt{staging.social\_flood\_signals}}

The naming choice was intentional. I call the rows ``signals'' rather
than ``events'' so nobody downstream confuses them with confirmed
flooding.

\begin{lstlisting}[language=SQL]
CREATE TABLE IF NOT EXISTS staging.social_flood_signals (
    id                      BIGSERIAL PRIMARY KEY,
    platform                TEXT NOT NULL,
    post_id                 TEXT NOT NULL,
    source_event_id         TEXT,
    created_at              TIMESTAMPTZ NOT NULL,
    ingested_at             TIMESTAMPTZ,
    text                    TEXT,
    language                TEXT,
    url                     TEXT,
    author_id_hash          TEXT,
    matched_keywords        TEXT[],
    excluded_keywords       TEXT[],
    flood_relevance_score   DOUBLE PRECISION,
    location_confidence     DOUBLE PRECISION,
    signal_confidence       DOUBLE PRECISION,
    place_name              TEXT,
    country                 TEXT,
    latitude                DOUBLE PRECISION,
    longitude               DOUBLE PRECISION,
    geometry                GEOMETRY(Point, 4326),
    h3_index                TEXT,
    raw_payload             JSONB,
    loaded_at               TIMESTAMPTZ NOT NULL DEFAULT NOW(),
    CONSTRAINT uq_social_flood_signals_platform_post
        UNIQUE (platform, post_id)
);
\end{lstlisting}

The schema migration is idempotent: every column is wrapped in
\texttt{ALTER TABLE \ldots\ ADD COLUMN IF NOT EXISTS}, and the unique
constraints are added inside a \texttt{DO} block that checks
\texttt{pg\_constraint} first. This was important because the
authoritative pipeline was already in production; I had to make sure my
schema deltas would apply cleanly to an existing Supabase instance, not
just to a fresh database.

I also added six task-specific indexes
(\texttt{platform}, \texttt{created\_at}, \texttt{country},
\texttt{h3\_index}, \texttt{signal\_confidence}, and a GIST index on
\texttt{geometry}) so the dashboard queries stay fast as the table grows.

\section{Ingestion: \texttt{ingestion/ingest\_bluesky.py}}

\subsection{Authentication Flow}

The single most important architectural detail of the ingester is which
host it talks to:

\begin{lstlisting}
PUBLIC_SEARCH_URL = "https://public.api.bsky.app/xrpc/app.bsky.feed.searchPosts"
AUTHED_SEARCH_URL = "https://bsky.social/xrpc/app.bsky.feed.searchPosts"
AUTH_URL          = "https://bsky.social/xrpc/com.atproto.server.createSession"
\end{lstlisting}

The choice is made on every request by a tiny helper:

\begin{lstlisting}
def _search_url(headers):
    return AUTHED_SEARCH_URL if headers else PUBLIC_SEARCH_URL
\end{lstlisting}

When \texttt{BLUESKY\_HANDLE} and \texttt{BLUESKY\_APP\_PASSWORD} are
configured, \texttt{\_auth\_headers()} calls
\texttt{com.atproto.server.createSession}, caches the returned
\texttt{accessJwt} in a process-level variable, and returns the
\texttt{Authorization: Bearer \ldots} header. The fetch helper then sends
the request to \texttt{bsky.social} instead of
\texttt{public.api.bsky.app}, because the public AppView does not accept
Bearer tokens.

\subsection{JWT Expiry Handling}

Access JWTs from Bluesky live roughly two hours. For a daily Airflow
schedule that is fine, but a long-running interactive session or a
backfill can outlive the token. Rather than refreshing eagerly on a
timer, I handle expiry reactively:

\begin{enumerate}[leftmargin=*]
  \item Send the search request with the cached JWT.
  \item If the response is 401 or 403 \emph{and} we were authenticated,
        clear the cached token, call \texttt{createSession} again, retry
        the request once.
  \item If the second attempt still fails with 401/403, raise an
        \texttt{HTTPError} whose message tells the operator exactly which
        environment variables to set.
\end{enumerate}

This keeps the happy path one HTTP call long, never refreshes
unnecessarily, and recovers from token expiry without operator
intervention.

\subsection{Bounded Search Window}

I added \texttt{BLUESKY\_LOOKBACK\_HOURS} to bound scheduled runs. When
set, \texttt{fetch\_posts} translates it into the \texttt{since}
timestamp Bluesky expects (RFC~3339 with a trailing~\texttt{Z}). An
explicit caller-supplied \texttt{since} always wins. The default in
production is~24 hours.

The reason is operational: without a bound, every Airflow run would
re-fetch the full~24-hour window plus everything older, and the de-dup
on \texttt{(platform, post\_id)} would silently make every run look
identical in cost. With the lookback, each run only pays for posts it
hasn't seen.

\subsection{Filtering Strategy}

A keyword-only filter is far too noisy. The classifier has four signal
groups feeding one decision:

\begin{itemize}[leftmargin=*]
  \item \textbf{Keywords} (\texttt{DEFAULT\_KEYWORDS}): trilingual
        English/French/Arabic flood vocabulary, both weak forms
        (``flood'', ``flooding'') and strong phrases (``flash flood'',
        ``levee breach'', \emph{inondation}, \emph{crue},
        \emph{fayadanat}). 30~entries.
  \item \textbf{Strong terms} (\texttt{STRONG\_FLOOD\_TERMS}): phrases
        unambiguous enough to admit a post on their own, without needing
        extra context.
  \item \textbf{Context terms} (\texttt{DEFAULT\_CONTEXT\_TERMS}):
        disaster vocabulary (``rain'', ``storm'', ``rescue'',
        ``evacuation'', ``levee'', ``road'', ``bridge'') in three
        languages. A weak keyword needs at least one of these to be
        admitted; otherwise the post is discarded.
  \item \textbf{Exclusions} (\texttt{DEFAULT\_EXCLUDED\_PHRASES} plus a
        compiled regex): obviously metaphorical phrases like ``flood my
        inbox'', ``flood of emails'', and a regex covering the
        ``flood~[the]~\textit{institution}'' political-protest family.
\end{itemize}

A post is kept only if:

\begin{lstlisting}
matched(any keyword)
AND NOT matched(any exclusion)
AND (matched(any strong term)
      OR matched(any context term)
      OR contains #flood or #flooding)
\end{lstlisting}

Every kept post records its \texttt{matched\_keywords},
\texttt{matched\_context\_terms}, \texttt{matched\_strong\_terms},
\texttt{excluded\_keywords}, \texttt{filter\_score},
\texttt{filter\_reason}, and a \texttt{source\_confidence} of~0.6 (the
post came from a public API with a stable identifier; that does not
make the claim true). All of this lives on the raw row so the
transformer never has to re-classify.

\subsection{Snapshot and Audit Trail}

Every run writes a JSON snapshot of the records it kept under
\texttt{data/raw/social\_media/bluesky/}, computes a SHA-256 of the
file, and writes one row into \texttt{raw.ingestion\_log} with the batch
id, status, row count, source URL actually used (authed or public),
file checksum, and timing. This mirrors what the five existing
ingestion modules do and lets me audit production runs against the
historical record.

\section{Transformation: \texttt{transformations/transform.py}}

\subsection{Normalization}

Each raw payload becomes one row in \texttt{staging.social\_flood\_signals}
via \texttt{\_normalize\_social\_media\_posts}. The function:

\begin{enumerate}[leftmargin=*]
  \item Pulls the post id, platform, created-at, text, language, URL,
        and author hash off the raw payload.
  \item Merges upstream exclusion markers with a fresh evaluation
        against the current rules (see self-healing below).
  \item Fills \texttt{country} and \texttt{place\_name} from text when
        the ingester didn't supply them (see geo extraction below).
  \item Computes three scores: flood relevance, location confidence,
        and the final signal confidence.
  \item Generates an H3 index if the row carries coordinates.
\end{enumerate}

\subsection{Scoring}

Three independent components feed one final score.

\paragraph{Flood relevance.} Uses the upstream \texttt{filter\_score}
as the base, then adds bonuses for additional matched keywords, context
terms, strong-term presence, and a small disaster-vocabulary check on
the text itself. If any current exclusion fires, relevance collapses
to~0.

\paragraph{Location confidence.} A simple five-tier function. I
deliberately kept the rule deterministic and short:

\begin{itemize}[leftmargin=*]
  \item~1.00 ---~latitude and longitude both present.
  \item~0.75 ---~place name and country both present.
  \item~0.50 ---~country only.
  \item~0.35 ---~place name only.
  \item~0.00 ---~no geographic clue.
\end{itemize}

\paragraph{Signal confidence.} A weighted combination:

\begin{lstlisting}
signal_confidence =
      0.50 * relevance
    + 0.25 * location_confidence
    + 0.15 * source_confidence
    + 0.10 * timestamp_score
\end{lstlisting}

clamped to~[0,~1]. Posts that fail relevance get a hard zero
regardless of the other components.

\subsection{Self-Healing Filter Re-Evaluation}

This is the design choice I'm most pleased with. When the transformer
normalises a raw row, it re-runs the \emph{current} exclusion rules
against the text. If a phrase matches now that didn't when the row was
first ingested, the row is flagged with the fresh exclusion. Then the
upsert helper:

\begin{itemize}[leftmargin=*]
  \item \texttt{INSERT \ldots\ ON CONFLICT DO UPDATE}s the valid signals.
  \item \texttt{DELETE}s any \texttt{(platform, post\_id)} whose
        \texttt{excluded\_keywords} just became non-empty.
\end{itemize}

The net effect: tightening the filter is a code change plus a
re-transform. No manual SQL surgery, no orphan rows lingering in
staging. The raw row stays in place, which is what I want for the
audit trail.

\subsection{Country and Place-Name Extraction}

The ingester sets \texttt{country} and \texttt{place\_name} to
\texttt{None} because Bluesky posts almost never carry structured place
metadata. To bring location coverage above zero I wrote a small
self-contained extractor at \texttt{transformations/social\_geo.py}.

The extractor runs five precision-first strategies, returning on the
first hit:

\begin{enumerate}[leftmargin=*]
  \item \textbf{Country adjective} (``French floods'' ->~France).
        70~adjectives, sorted longest-first so ``Sri Lankan'' beats
        ``Lankan''.
  \item \textbf{Explicit country name} (``\ldots across Spain
        \ldots'' ->~Spain). About~110~countries, expressed as compiled
        case-insensitive regexes so ``UK'', ``U.K.'', ``United Kingdom''
        all map to one canonical form.
  \item \textbf{\texttt{<City>, <ST>} pattern} ->~United States with a
        full \texttt{place\_name} like ``Baton Rouge, Louisiana''. The
        city group is a title-cased 1--4-word sequence, the state group
        is a two-letter all-caps code from a closed list of the
        50~states plus DC.
  \item \textbf{Full US state name} ->~United States with the state
        name as \texttt{place\_name}.
  \item \textbf{US timezone abbreviation} (\texttt{AKDT},
        \texttt{EDT}, \ldots) ->~United States. Specifically targets
        NWS weather alerts which routinely embed the local timezone
        code.
\end{enumerate}

The module is zero-dependency (only \texttt{re}) and never overwrites a
value the ingester already supplied. If a future Mastodon ingester
delivers a structured place field, the extractor stays out of the way.

The whole module is under~300~lines, runs in microseconds, and is
covered by 28~tests anchored on real Bluesky samples I observed during
the cutover. The negative tests are as important as the positive
ones~---~``OR'' as an English conjunction must not be misread as
Oregon, and a political metaphor must not produce a misleading country.

\section{Marts}

Three views are rebuilt by \texttt{transformations/marts.py} on every
DAG run. The views are dropped and recreated rather than refreshed
in place because PostgreSQL's \texttt{CREATE OR REPLACE VIEW} only
succeeds if the column list is byte-identical, and I wanted to be able
to evolve them without manual migrations.

\subsection{\texttt{marts.social\_flood\_signals}}

A flat projection over staging with the geometry rendered as GeoJSON
for the FastAPI layer.

\subsection{\texttt{marts.flood\_events\_with\_social\_signals}}

A LEFT JOIN of the deduplicated event view onto staging signals,
aggregated per event. A signal corroborates an event if its
\texttt{created\_at} falls inside [\texttt{date\_start - 1 day},
\texttt{date\_end + 3 days}] and either the H3 cells match or the
country names match case-insensitively. The output adds
\texttt{social\_signal\_count},
\texttt{first\_social\_signal\_at},
\texttt{latest\_social\_signal\_at},
\texttt{avg\_social\_signal\_confidence}, and a deduplicated
\texttt{social\_platforms} array.

\subsection{\texttt{marts.social\_signals\_by\_country\_day}}

A daily, per-country rollup designed for monitoring dashboards. Each
row has the signal count, the average confidence, the count of
high-confidence (\texttt{>= 0.5}) signals, the set of platforms, and
the set of languages. \texttt{country = NULL} is mapped to the literal
string \texttt{'unknown'} so the consumer never has to handle
\texttt{NULL} grouping behavior. This is the view I expect to wire
into the next dashboard iteration.

\section{API: \texttt{api/main.py}}

The existing FastAPI service was extended with four social endpoints,
each of which guards its mart with a \texttt{503} if the view does not
yet exist (so the API is honest about pipeline state rather than
returning misleading empty arrays):

\begin{itemize}[leftmargin=*]
  \item \texttt{GET /social-signals} ---~paginated signal feed with
        \texttt{platform}, \texttt{country}, and
        \texttt{min\_confidence} filters.
  \item \texttt{GET /flood-events/with-social-signals} ---~events
        enriched with corroborating signal counts; filterable by
        country and \texttt{min\_social\_signals}.
  \item \texttt{GET /analytics/social-signals/by-platform} ---~one row
        per platform with totals and average confidence.
  \item \texttt{GET /analytics/social-signals/by-country-day} ---~the
        country/day rollup with optional country, start/end date, and
        confidence-floor filters.
\end{itemize}

\section{Orchestration: \texttt{dags/flood\_event\_pipeline\_dag.py}}

The Airflow DAG was extended with one new task. The wiring keeps
social ingestion parallel with the five authoritative ingestions:

\begin{lstlisting}
apply_schema
  -> validate_endpoints
       -> ingest_dartmouth ------+
       -> ingest_glofas ---------+
       -> ingest_copernicus_ems -+--> transform_raw_to_staging
       -> ingest_emdat ---------+        -> refresh_marts
       -> ingest_reliefweb ----+              -> dbt_run
       -> ingest_bluesky ----+                      -> data_quality
\end{lstlisting}

\texttt{ingest\_bluesky} is no-op unless
\texttt{ENABLE\_SOCIAL\_MEDIA\_INGESTION} is truthy in the environment,
so toggling social on or off is a single env-var change with no DAG
edits. The \texttt{data\_quality} task uses \texttt{trigger\_rule =
all\_done} so an upstream ingestion failure never silences the report.

\section{Data Quality}

\texttt{validation/data\_quality.py} runs eight social-signal checks
alongside the existing flood-event checks. Every check is a simple
\texttt{SELECT COUNT(*) AS bad FROM \ldots} returning the number of
failing rows; the report renders them as a Markdown table at
\texttt{data/logs/data\_quality\_report.md}.

\begin{itemize}[leftmargin=*]
  \item \texttt{social\_duplicate\_platform\_post\_ids}
  \item \texttt{social\_missing\_created\_at}
  \item \texttt{social\_missing\_platform\_or\_post\_id}
  \item \texttt{social\_invalid\_latitude\_or\_longitude}
  \item \texttt{social\_missing\_h3\_with\_coords}
  \item \texttt{social\_confidence\_out\_of\_range}
  \item \texttt{social\_relevance\_without\_keywords}
        \emph{(integrity: future refactor must not let a row land in
        staging without keywords)}
  \item \texttt{social\_orphan\_staging\_signals}
        \emph{(traceability: every staging row must still point at an
        existing raw row)}
\end{itemize}

The report also includes a per-platform summary table mirroring the
per-source table for authoritative events.

\section{Problems Faced and How I Solved Them}

\subsection{The 401/403 Wall on Production Cutover}

\textbf{Symptom.} Local tests all green, schema applied to Supabase,
ingester refused to fetch anything: every call returned~401 or~403.
Filling \texttt{BLUESKY\_HANDLE} and \texttt{BLUESKY\_APP\_PASSWORD} in
\texttt{.env} did not change a thing.

\textbf{Root cause.} The ingester was hard-coded to
\texttt{public.api.bsky.app}, Bluesky's \emph{unauthenticated} AppView.
That host does not accept Bearer tokens. Sending one is at best
ignored and at worst rejected as malformed. So the credentials I had
configured were being silently dropped before the request even left my
machine.

\textbf{Fix.} I separated the two hosts into two named constants and
introduced \texttt{\_search\_url(headers)}, which picks the right host
based on whether the request carries auth headers. Authenticated calls
now go to \texttt{bsky.social/xrpc/app.bsky.feed.searchPosts}~---~the
PDS, which proxies authenticated reads to the AppView and accepts the
JWT returned by \texttt{createSession}. I also widened the error
message so a future operator sees the actionable
``set \texttt{BLUESKY\_HANDLE} and \texttt{BLUESKY\_APP\_PASSWORD}''
hint instead of a generic ``HTTP~401''.

\textbf{Aftermath.} The \texttt{raw.ingestion\_log} still carries the
historical failure row with \texttt{source\_url} pointing at the public
AppView. I kept it deliberately as evidence that the fix was effective:
every subsequent row records \texttt{bsky.social} as the actual host.

\subsection{The Leading @ on \texttt{BLUESKY\_HANDLE}}

\textbf{Symptom.} After fixing the host issue, my next attempt still
failed. The \texttt{createSession} call returned HTTP~400 with no
\texttt{accessJwt}.

\textbf{Root cause.} The Bluesky web UI displays handles with a leading
\texttt{@} (``\texttt{@user.bsky.social}''). I copied mine verbatim into
\texttt{.env}. The XRPC endpoint expects the bare handle and rejects
the \texttt{@}-prefixed form as malformed input.

\textbf{Fix.} \texttt{\_bluesky\_credentials()} now strips a leading
\texttt{@} and trims surrounding whitespace before returning the
identifier. I covered both the normalization and the
``\texttt{@}-only handle'' degenerate case with tests. The next operator
who copies a handle out of the UI never has to debug this again.

\subsection{Pytest \texttt{tmp\_path} on Windows}

\textbf{Symptom.} \texttt{pytest} refused to run on my Windows machine
with \texttt{PermissionError: [WinError 5] Access is denied} on
\texttt{C:\textbackslash{}Users\textbackslash{}\ldots\textbackslash{}Temp\textbackslash{}pytest-of-\ldots}.

\textbf{Root cause.} An old \texttt{pytest-of-\ldots} directory left by
a different Windows user had locked permissions; pytest could neither
delete nor reuse it.

\textbf{Fix.} I run the suite with an explicit, writable
\texttt{--basetemp=C:\textbackslash{}Temp\textbackslash{}pytest\_flood}
and \texttt{-p no:cacheprovider} to bypass the \texttt{.pytest\_cache}
folder my IDE had also locked. This is a workspace-specific
workaround, not a code change.

\subsection{Filter Admitting Political Metaphors}

\textbf{Symptom.} The first 50 production rows looked good on first
glance, but a manual inspection of the lowest-confidence signals showed
several political-protest tweets: ``we have to flood Congress'', ``flood
the zone with shit'', ``We should be flooding the streets''.

\textbf{Root cause.} The old exclusion list covered the obvious
metaphors (``flood my inbox'') but not the political-action family. The
word ``Congress'' and ``streets'' are not in the disaster-context list
but neither was anything else strong enough to block them, and the
weak keyword ``flood'' was enough to admit them when paired with the
generic alert vocabulary the same posts often carried.

\textbf{Fix.} I added fourteen new literal exclusion phrases plus a
compiled regex that matches the family generically:

\begin{lstlisting}
flood(?:ing|ed)? \s+ (?:the\s+|up\s+the\s+)?
   (?: zone | portal | inbox | congress | senate |
       white\s+house | polls? | courts? | streets? |
       airwaves | system | government | hill | capitol |
       comments? | timeline | mentions? | feed |
       chat | chats | dms? | replies? )
\end{lstlisting}

When the regex fires, \texttt{\_excluded\_keywords} records a synthetic
\texttt{political\_metaphor:flood\_the\_<institution>} marker on the
row, so the audit trail explains exactly which rule rejected the post.

To make the fix retroactive without touching production by hand, the
transformer re-evaluates the live exclusion rules at every transform
pass (see the self-healing design earlier). Running the transformer
once after the change dropped four political metaphors out of the
existing~50 staging rows automatically.

\subsection{One Hundred Percent of Signals Tagged \texttt{country = NULL}}

\textbf{Symptom.} The new \texttt{marts.social\_signals\_by\_country\_day}
view collapsed all~50 production signals into a single
\texttt{country = 'unknown'} bucket, which made it useless for
monitoring.

\textbf{Root cause.} Bluesky posts almost never carry structured place
fields, so the ingester always sets \texttt{country} and
\texttt{place\_name} to \texttt{None}.

\textbf{Fix.} The new \texttt{transformations/social\_geo.py} extractor
described earlier. The change is fully backward-compatible: it only
runs when \texttt{country} or \texttt{place\_name} is missing, and it
returns \texttt{None} when no strategy is confident. After
re-transforming the same~50 raw rows, country coverage went from~0\%
to~61\% and place-name coverage from~0\% to~50\%, with the United
States, Spain, Pakistan, and Turkey appearing as distinct buckets in
the country/day mart.

\subsection{Local Python Environment Without \texttt{h3}}

\textbf{Symptom.} Trying to install the project's full dependency set
into a single local Python environment failed: \texttt{h3} requires a
C~toolchain on Windows that isn't installed, and \texttt{pandas}
attempted to build from source for a Python version with no prebuilt
wheel.

\textbf{Fix.} Production code imports \texttt{h3} normally and falls
back to \texttt{None} only if the import fails. The transformer logs
a warning and leaves \texttt{h3\_index} \texttt{NULL} when the package
is missing, so the pipeline keeps producing valid staging rows in
locally-degraded environments. The test suite includes a tiny H3 shim
that activates only when the real package is unavailable, so the
transformation tests still exercise the \texttt{h3\_index} code path.
Production Airflow runs inside the Docker image where the real
\texttt{h3} package is baked in.

\subsection{Friend's Supabase Project Visibility}

\textbf{Symptom.} The data pipeline writes succeeded but the project
owner couldn't see the new tables in their own Supabase dashboard.

\textbf{Root cause.} The \texttt{DATABASE\_URL} in \texttt{.env}
points to a Supabase project owned by a different account
(\texttt{project\_ref = grakjgjmuusbrbjlgvgg}). Logging into Supabase
with a different account doesn't show that project.

\textbf{Fix.} Documented the project reference and the dashboard URL
in the verification script and the runbook so the next operator either
asks the project owner for a Team invite or signs in with the owning
account.

\subsection{PowerShell Quote Mangling in One-Liner Verification}

\textbf{Symptom.} A throwaway \texttt{python -c "\ldots"} command for
inspecting the database failed with \texttt{SyntaxError: unterminated
string literal} because PowerShell pre-processed the embedded SQL
string.

\textbf{Fix.} Moved verification SQL into a small dedicated script
(\texttt{scripts/\_verify\_social\_cutover.py}) that I can re-run as a
single command. It also doubles as a sanity-check tool for the next
deploy.

\section{Live Production Results}

After all fourteen steps, against the real Supabase project, on a
72-hour Bluesky lookback window with bounded sizes
(\texttt{max\_posts=50}, \texttt{per\_query\_limit=25}):

\begin{center}
\small
\begin{tabular}{lrr}
\toprule
Metric & After cutover only & After filter + geo \\
\midrule
Raw rows ingested                & 50    & 50    \\
Staging signals                  & 50    & 46    \\
Average \texttt{signal\_confidence}     & 0.651 & 0.766 \\
Rows with \texttt{country}              & 0     & 28    \\
Rows with \texttt{place\_name}          & 0     & 23    \\
Distinct countries in staging    & 1     & 5     \\
Failing social DQ checks         & 0 / 8 & 0 / 8 \\
\bottomrule
\end{tabular}
\end{center}

Cross-source corroboration is not yet a measurable signal because the
first 50 social rows fell outside the time/H3/country windows of every
authoritative event currently in \texttt{staging.flood\_events}.
That is expected at the scale of a single small ingestion run; it will
become measurable once daily Airflow runs accumulate hundreds of
signals per week.

\section{Testing}

138 automated tests cover the social pipeline, organised as:

\begin{itemize}[leftmargin=*]
  \item \texttt{tests/test\_ingest\_bluesky.py} (29 tests)~---~public
        vs.\ authed host selection, JWT-refresh retry, lookback window,
        explicit-\texttt{since} override, credential normalization
        (including the leading-\texttt{@} case),
        political-metaphor exclusion (8~parametrized rejects,
        4~parametrized keeps), regex sanity, snapshot writing,
        ingestion-log success and failure paths.
  \item \texttt{tests/test\_transform.py} (52 tests)~---~scalar
        coercers, per-source normalisers for the five authoritative
        sources, social normaliser including the new geo enrichment
        and the don't-overwrite invariant, self-healing exclusion
        re-evaluation, upsert filtering and the DELETE pass for newly
        excluded rows.
  \item \texttt{tests/test\_social\_geo.py} (28~tests)~---~real-sample
        regression cases, country adjectives, country names,
        \texttt{<City>, <ST>} patterns, full state names, US timezone
        codes, plus negative cases for empty input, ambiguous
        two-letter words, and political metaphors that must not
        produce country tags.
  \item \texttt{tests/test\_db\_client.py} (16~tests)~---~including
        the social insert helper, NaN/Inf cleaning, idempotent
        upserting, transient-pooler-drop retry.
  \item \texttt{tests/test\_ingestion\_common.py}
        (10~tests)~---~shared helpers used by every ingester.
\end{itemize}

The suite runs in roughly one second on my machine. None of these
tests touch the network or the live database; all external dependencies
are mocked.

\section{Future Enhancements}

These are the changes I'd prioritise next, in rough order of value per
effort.

\begin{enumerate}[leftmargin=*]
  \item \textbf{Offline gazetteer for non-US cities.} The current geo
        extractor is precision-first and recognises explicit country
        mentions plus US states and timezones. Names of major cities
        outside the US (Lagos, Karachi, Manila, S\~ao Paulo) are still
        invisible to it. A small static gazetteer (e.g.\ GeoNames
        \texttt{cities500} restricted to >\,1\,M population) would lift
        recall without bringing in a heavyweight NLP dependency.
  \item \textbf{Cross-platform corroboration.} Add a second platform
        (Mastodon search is the easiest because of its public API and
        absence of strict rate limits). Strong corroboration should
        require independent signals from at least two platforms, not
        just multiple posts on Bluesky.
  \item \textbf{Social-incident clustering.} Group signals by
        country / H3 / time window into candidate incidents and
        compute a cluster confidence score that rewards quantity,
        independence (distinct authors), and source diversity.
  \item \textbf{Alerting hook.} When
        \texttt{marts.social\_signals\_by\_country\_day.high\_confidence\_count}
        crosses a threshold for a country with no recent matching
        event in \texttt{marts.flood\_events\_unique}, fire a
        notification. This is the operational reason the country/day
        mart exists in the first place.
  \item \textbf{Bounded AT~Protocol firehose.} If \texttt{searchPosts}
        coverage proves insufficient for early-warning, the natural
        next step is a bounded firehose client with a strict
        budget~---~not a long-running daemon, but a windowed reader
        that runs alongside the search-based ingester.
  \item \textbf{Dashboard panel.} Surface the new country/day rollup
        and the per-event signal corroboration in the dashboard so the
        evidence layer is visible alongside authoritative events.
\end{enumerate}

\section{Conclusion}

What began as a request to ``add social posts to the pipeline'' turned
into a 14-step exercise in keeping a new evidence layer separate from
authoritative data, handling Bluesky's two-tier authentication model
correctly, building a filter that is both explainable and tightenable
without producing orphan rows, and recovering geographic context from
free text without either an external service or a heavy dependency.
The system is now live, instrumented, tested, and operationally
honest: every score and every rejection is justified by data on the
row itself, every run is auditable through
\texttt{raw.ingestion\_log}, and every dashboard query is answered by
a mart that explicitly tells the API when it is not ready.

\end{document}
